\begin{claim}{D.2.22}{Proved}{The tenth-stage target: exponent, integral Hurewicz kernel, and eta composites}
Keep $B=K(\FF_2,2)$, $Y=\Sigma^2B$, $T=\Omega^2Y$, and the
three-connected fiber $F$ of \textup{D.2.18--D.2.20}.
Let $\rho\in\pi_6T$ be the order-four element of
\textup{D.2.18}. Then
\[
 4\pi_9T=0,\qquad
 \ker\bigl(h_{\ZZ}:\pi_9T\longrightarrow H_9(T;\ZZ)\bigr)
   =\ZZ/2\{\rho\nu_6\},\qquad \rho\nu_6\ne0,
 \tag{D.2.22a1}
\]
and
\[
 \pi_8T\circ\eta_8=0.
 \tag{D.2.22a2}
\]
The group $\pi_9T$ is either $(\ZZ/2)^4$ or
$\ZZ/4\oplus(\ZZ/2)^2$. The calculation retains the unresolved
common even-nu attachment coefficient that distinguishes these
two alternatives. For the bottom class $\chi\in\pi_4F$ used in
\textup{D.2.21}, one has the precise restriction
\[
 \lambda_1(\chi,\chi)\in\{0,\rho\nu_6\}.
 \tag{D.2.22a18}
\]
Its value is not selected. The criterion of \textup{D.2.21}
therefore reduces to vanishing of this actual square. No extension
of the combined map through the tenth skeleton is asserted.
\end{claim}
\begin{proof}
All displayed homology operations have coefficients $\FF_2$ unless
integral coefficients are indicated. The spaces have finite
two-primary homotopy groups in the degrees used, by the Serre-class
argument already applied to $Y$. Thus the two-local component
calculation below computes the actual groups in the statement.

\medskip\noindent
\textbf{The complete next fiber homology table.}
Continue the label notation of \textup{D.2.20}. Put
\[
 \begin{gathered}
 A=a^2,\quad z=e+ab,\quad u_F=\xi a,\quad
 k=\{a,b\},\quad R=\{a,c\},\\
 d_6=b^2,\quad e_6=b^2+k,\quad f_6=f+ac+a^3.
 \end{gathered}
 \tag{D.2.22a3}
\]
Here the unadorned $d,e$ remain the degree-five label classes,
dual to $w_B$ and $u_Bv_B$; they are not $d_6,e_6$.
The previously computed classes are
\[
\begin{gathered}
 P=g+ae+bc+a^2b,\qquad Q=h+ad,\\
 W=a^4,\quad C=c^2,\quad L=l+bd,\quad
 M=j+ax+be+ab^2,\\
 J_1=\{a,d\},\quad J_2=\{a,e\}+ak,\quad
 J_3=\{b,c\}+ak.
\end{gathered}
\]
Write $p,q,r,s$ for the new degree-ten labels dual to
$u_B^5,u_B^2v_B^2,u_Bv_Bw_B,w_B^2$.
The free double-loop basis has twelve indecomposables in degree ten:
\[
\begin{gathered}
 p,q,r,s,\\
 \{a,g\},\{a,h\},\{b,f\},\{b,x\},
 \{c,d\},\{c,e\},\{b,k\},\{a,R\}.
\end{gathered}
\]
The restricted identities identify $\{a,\xi b\}$ with $\{b,k\}$
and $\{c,\xi a\}$ with $\{a,R\}$.
The thirty products are
\[
\begin{gathered}
 ai,aj,al,a\{a,d\},a\{a,e\},a\{b,c\},\\
 bg,bh,b\xi b,bR,cf,cx,ck,\\
 d^2,de,du_F,e^2,eu_F,u_F^2,\\
 a^2f,a^2x,a^2k,abd,abe,abu_F,\\
 ac^2,b^2c,a^3c,a^2b^2,a^5.
\end{gathered}
\]
Thus $\dim H_{10}T=42$. The base and mixed summands of
$H_{10}(B\times F)$ have dimensions $4+2+4+3+4+7=24$,
so $\dim H_{10}F=18$.

A complete fiber basis consists of
\[
\begin{gathered}
 D_A=Ad_6,\quad E_A=Ae_6,\quad F_A=Af_6,
 \quad Z_2=z^2,\quad ZU=zu_F,\quad U_2=u_F^2,
 \quad D_2=d^2,\\
 Q_{10}=q+cx+aM+bP+abz,\\
 P_p=p+ai+cf+a^2f+ac^2+a^3c+a^5,
 \qquad P_r=r+al+bh+ed,
\end{gathered}
\tag{D.2.22a4}
\]
and the eight classes
\[
\begin{aligned}
 B_1&=\{a,g\}+a\{a,e\}+ck+bR+Ak,&
 B_2&=\{a,h\}+a\{a,d\},\\
 B_3&=\{b,f\}+ck+a\{b,c\}+Ak,&
 B_4&=\{b,x\},\\
 B_5&=\{c,d\}+a\{a,d\},&
 B_6&=\{c,e\}+bR+a\{b,c\}+a\{a,e\}+Ak,\\
 B_7&=\{b,k\},&B_8&=\{a,R\}.
\end{aligned}
\tag{D.2.22a5}
\]
The classes $D_2,P_p,P_r,B_1,\ldots,B_8$ are primitive,
and $\overline\Delta Q_{10}=z\otimes z$.
For example
\[
\begin{aligned}
 \overline\Delta\{a,g\}
 &=a\otimes\{a,e\}+\{a,e\}\otimes a
   +k\otimes c+c\otimes k+b\otimes R+R\otimes b,\\
 \overline\Delta\{b,f\}
 &=k\otimes c+c\otimes k
   +a\otimes\{b,c\}+\{b,c\}\otimes a,\\
 \overline\Delta\{c,e\}
 &=b\otimes R+R\otimes b
   +a\otimes(\{b,c\}+\{a,e\})\\
 &\quad +(\{b,c\}+\{a,e\})\otimes a
   +A\otimes k+k\otimes A.
\end{aligned}
\]
The other bracket diagonals have respectively the two terms canceled
by the displayed product correction; $\{b,x\}$ is primitive because
$\overline\Delta x=b\otimes b$ and $\{b,b\}=0$.
The polynomial diagonal of $Q_{10}$ follows by expanding it as
$q+cx+aj+a^2x+abe+a^2b^2+bg+b^2c$.
All these classes have zero base image; there is no primitive in
$H_{10}B$. Their distinct leading monomials prove independence,
and the dimension count proves completeness.

Write $K_\Sigma=K_2+K_3+K_4$ and
$Z_0=P_m+Z_A+P_n$ in the degree-nine notation of
\textup{D.2.20}. The entire operation table is
\[
\begin{array}{c|ccc}
 &\beta&\operatorname{Sq}_*^2&\operatorname{Sq}_*^4\\\hline
 D_A,E_A,F_A,Z_2&0&W&0\\
 ZU&U_A&0&0\\
 U_2&0&0&0\\
 D_2&0&0&d_6\\
 Q_{10}&Z_0&M+C+W&f_6\\
 P_p&0&W&f_6\\
 P_r&0&0&0\\
 B_1&K_1&J_2&0\\
 B_2&0&J_2&0\\
 B_3&K_1&0&0\\
 B_4&K_\Sigma&J_1+J_2+J_3&d_6+e_6\\
 B_5&0&J_1+J_3&d_6+e_6\\
 B_6&0&J_2&0\\
 B_7&K_5&0&0\\
 B_8&0&0&0
\end{array}
\tag{D.2.22a6}
\]
It follows directly from Cartan and the CLM Browder identities
\cite[Chapter III, Theorems 1.2--1.4]{I_CohenLadaMay1976}.
For the new labels, the triples $(\beta,\operatorname{Sq}_*^2,
\operatorname{Sq}_*^4)$ are $(0,0,f)$ for $p$,
$(m+n,j,f)$ for $q$, $(0,0,0)$ for $r$, and $(t,l,x)$ for $s$.
In particular $\beta Q_{10}=Z_0$, an independent check on the
essential second-Bockstein correction in \textup{D.2.20}.
The coproducts and twenty-one operation values on the seven
polynomial classes are also independently checked by
\texttt{code/FGH\_H10\_poly\_check.py} and the root's separate
polynomial certificate; the bracket calculations above do not
depend on those certificates.

\medskip\noindent
\textbf{The order-eight class and both third Bocksteins.}
The universal-coefficient theorem, the known $H_8F$, and the
rank-five first Bockstein in \textup{(D.2.22a6)} give six cyclic
summands in $H_9(F;\ZZ)$, five of order two. The remaining
higher-exponent class is represented by $N$.
We determine its exact order, including the second order-eight
boundary that will be essential below.

Use the integral chain summand of $B$ with $c$ in degree four,
$d$ in degree five, and $\partial d=4c$. In the arity-two
configuration-circle complex put
\[
 z_c=e_1(cc),\quad v=e_0(cd-dc),\quad
 q_0=e_0(dd),\quad h=e_1(cd),\quad h'=e_1(dc).
\]
The Koszul transposition signs give the exact formulas
\[
\begin{gathered}
 \partial q_0=4v,\qquad
 \partial h=-v-4z_c,\qquad
 \partial h'=v-4z_c,\\
 \partial e_1(dd)=-2q_0-4h+4h'.
\end{gathered}
\tag{D.2.22a7}
\]
The degree-nine cycles have basis $v,z_c$, and the boundary
columns are $(4,0),(-1,-4),(1,-4)$, with Smith invariants
$1,8$. Thus the class of $z_c$ has exact order eight.
Both mixed-label circle chains are needed: using only one
would miss the relation $\partial(h+h')=-8z_c$.

This chain calculation is a retract of the actual arity-two
summand. Indeed a bounded-below free integral chain complex
has a Smith decomposition, so $\widetilde C_*(B;\ZZ)$ retracts
as a chain complex onto the displayed Moore-four summand.
Its tensor square is equivariant for the signed transposition.
The configuration chains are free over $\ZZ[\Sigma_2]$;
the equivariant shuffle equivalence therefore identifies the
derived tensor with the chains of the actual filtration quotient.
The integral stable splitting of Snaith
\cite[Theorem 1.1 and Section 2]{Snaith1974} makes this a
homology retract of $T$. No topological Moore-space retract
of $B$ is assumed.

The integral lift of $N$ has weight-two component $\xi c$;
the other terms $aR$ and $a^2\xi a$ have order dividing two.
Thus it has order eight in $H_9T$. In the actual product splitting
$T\simeq B\times F$, all summands outside $H_9F$ have exponent
at most four. This includes the mixed Kunneth summands, notably
$H_4B\otimes H_5F\cong\ZZ/4$; the base summand itself has
exponent two by the integral Eilenberg--Mac Lane calculation
\cite[\S3.2]{I_Clement2002}. Projection therefore retains order eight in the
fiber. Consequently
\[
 H_9(F;\ZZ)\cong\ZZ/8\oplus(\ZZ/2)^5.
 \tag{D.2.22a8}
\]

Choose the integral labels so that $e=d+y$, $\partial y=0$,
and $\partial x=2y$, allowing odd-unit changes. The same
Moore-four retraction sends both $d$ and $e$ to $d$.
Hence the weight-two parts $\{c,d\}$ of $B_5$ and
$\{c,e\}$ of $B_6$ both project to $h+h'$.
Their higher-weight corrections disappear. This proves
\[
 \beta_3[B_5]=[N],\qquad \beta_3[B_6]=[N].
 \tag{D.2.22a9}
\]
For an exhaustive check, the first-Bockstein cycle basis in
degree ten is
\[
 D_A,E_A,F_A,Z_2,U_2,D_2,P_p,P_r,
 B_2,B_5,B_6,B_8,B_1+B_3.
\]
Only $B_5,B_6$ project to $h+h'$; the others project either
to zero or to $q_0$. The latter is already a first-Bockstein
boundary by \textup{(D.2.22a7)}. Thus these are exactly the
two nonzero third differentials into $N$. Projection also
shows that $N$ cannot be hit earlier and that these two
classes survive to the third page.

\medskip\noindent
\textbf{The four-connected cover and the complete low Koszul check.}
Let $E=F\langle5\rangle$ be the fiber of the first Postnikov
map $p:F\to K(\FF_2,4)$, representing $A^*$.
Serre's polynomial description \cite[Section 2]{I_Serre1953}
gives the following complete base-generator list through degree eleven:
\[
\begin{gathered}
 x_4=x,\ x_5=\operatorname{Sq}^1x,\ x_6=\operatorname{Sq}^2x,
 \ x_7=\operatorname{Sq}^3x,\ y_7=\operatorname{Sq}^2\operatorname{Sq}^1x,\\
 x_8=\operatorname{Sq}^3\operatorname{Sq}^1x,
 \ x_9=\operatorname{Sq}^4\operatorname{Sq}^1x,
 \ x_{10}=\operatorname{Sq}^4\operatorname{Sq}^2x,\\
 x_{11}=\operatorname{Sq}^5\operatorname{Sq}^2x,
 \ y_{11}=\operatorname{Sq}^4\operatorname{Sq}^2\operatorname{Sq}^1x.
\end{gathered}
\]
Their images through degree ten are
\[
\begin{array}{c|l}
 x_4&A^*\\
 x_5&z^*\\
 x_6&d_6^*+e_6^*+f_6^*\\
 x_7&V^*+P^*\\
 y_7&P^*+Q^*\\
 x_8&L^*+M^*\\
 x_9&P_n^*+P_o^*\\
 x_{10}&D_2^*+Q_{10}^*+P_p^*.
\end{array}
\tag{D.2.22a10}
\]
For degree eleven let $v_{11}$ be the label dual to
$u_B\operatorname{Sq}^4w_B$, and retain the degree-nine label
$t$ dual to $\operatorname{Sq}^4w_B$. The primitive fiber
classes $\xi d$ and $v_{11}+at$ satisfy
\[
\begin{array}{c|cc}
 &\beta&\operatorname{Sq}_*^4\\\hline
 \xi d&D_2&V\\
 v_{11}+at&0&Q.
\end{array}
\tag{D.2.22a11}
\]
The top CLM formulas give the first row: the only fourth-square
term is $Q^4b=\xi b$. Label Cartan gives the second row.
There is no primitive in $H_{11}B$, so both lie in the fiber.
They detect $x_{11},y_{11}$ as the identity matrix.
All base products through degree eleven are independently detected
by $C,Z_A,D_A,E_A,F_A,Q_{10}$ and the products
$AH_7F,zH_6F$. Thus $p^*$ is injective through degree eleven.

Injectivity alone would not exclude mixed Serre differentials.
Instead use the convergent cohomological Eilenberg--Moore sequence
\cite[p.334]{FGHMay1968} and its polynomial Koszul complex.
A Koszul generator for a base generator of degree $j$ has total
degree $j-1$. Length three starts in degree $3+4+5=12$;
length two with a positive $H^*F$ coefficient starts in degree
eleven. Hence the entire positive-length calculation through ten
has scalar length-two coefficients and length-one coefficients
of degree at most seven. The exact ranks are
\[
\begin{array}{c|ccc}
 \text{total degree}&\dim K_1&\operatorname{rank}(K_1\to K_0)
 &\operatorname{rank}(K_2\to K_1)\\\hline
 7&2&2&0\\
 8&4&3&1\\
 9&7&6&1\\
 10&13&10&3.
\end{array}
\tag{D.2.22a12}
\]
In degrees three through six the scalar images are independent.
For degree seven the images are $C^*,L^*+M^*$.
For degree eight, $A^*H^5F$ has rank two and $p^*x_9$
adds one; the single relation is the Koszul relation for
$x_4,x_5$. For degree nine, $A^*H^6F$ and $z^*H^5F$
have disjoint ranks three and two, and $p^*x_{10}$ adds one
rank detected on primitive $D_2$; the single relation is
the one for $x_4,x_6$.

For degree ten, $A^*H^7F$ and $z^*H^6F$ have disjoint ranks
four and three. The product $u_F^*p^*x_6$ adds one, while
$z^*p^*x_6$ is already in the latter rank-three space.
The two primitive tests in \textup{(D.2.22a11)} add two more
ranks. The three relations are those for
$(x_4,x_7),(x_4,y_7),(x_5,x_6)$.
All product ranks follow by evaluation on products of the
primitive homology factors; the equal-factor product $z^{*2}$
is detected on $Q_{10}$, not on $z^2$.
The scalar length-two wedges in degrees seven through ten are
\[
 t_3t_4;\quad t_3t_5;\quad
 t_3t_6,t_3t'_6,t_4t_5;\quad
 t_3t_7,t_4t_6,t_4t'_6.
\]
Their differentials are independent. This proves that positive
Tor vanishes through total degree ten. Higher base generators
have Koszul degree at least eleven and cannot contribute.
There is consequently no incoming higher Eilenberg--Moore
differential in this range, and restriction gives
\[
 H^{\le10}E\cong
 \bigl(H^*F/(p^*\widetilde H^*K(\FF_2,4))H^*F\bigr)^{\le10}.
 \tag{D.2.22a13}
\]

\medskip\noindent
\textbf{The cover table used by the component calculation.}
Dually $H_iE\to H_iF$ is injective through degree ten, with
the following complete image basis:
\[
\begin{array}{c|l}
5&u_F\\
6&D=d_6+e_6,\quad E_6=d_6+f_6\\
7&S=V+P+Q,\quad R\\
8&W,\quad D_8=L+M,\quad J_1,J_2,J_3\\
9&N,K_1,K_2,K_3,K_4,K_5,P_m,S_9=P_n+P_o\\
10&Z_2,U_2,T_{10}=D_2+P_p,P_r,B_1,\ldots,B_8.
\end{array}
\tag{D.2.22a14}
\]
Their dimensions are $1,2,2,5,8,12$, and every displayed
class is primitive. Below degree nine the nonzero first,
second, and fourth squares are precisely
\[
 \beta S=D+E_6,\quad \operatorname{Sq}_*^2S=u_F,
 \quad \operatorname{Sq}_*^2D_8=E_6,
 \quad \operatorname{Sq}_*^2J_1=\operatorname{Sq}_*^2J_3=D,
 \quad\beta J_3=R.
\]
In degree nine the table is
\[
\begin{array}{c|ccc}
 &\beta&\operatorname{Sq}_*^2&\operatorname{Sq}_*^4\\\hline
 N&0&R&u_F\\
 K_1&0&0&0\\
 K_2&J_1+J_2&R&0\\
 K_3&J_1&0&0\\
 K_4&J_2&0&0\\
 K_5&0&0&0\\
 P_m&W&0&0\\
 S_9&D_8&0&0
\end{array}
\tag{D.2.22a15}
\]
and in degree ten it is
\[
\begin{array}{c|ccc}
 &\beta&\operatorname{Sq}_*^2&\operatorname{Sq}_*^4\\\hline
 Z_2&0&W&0\\
 U_2&0&0&0\\
 T_{10}&0&W&E_6\\
 P_r&0&0&0\\
 B_1&K_1&J_2&0\\
 B_2&0&J_2&0\\
 B_3&K_1&0&0\\
 B_4&K_\Sigma&J_1+J_2+J_3&D\\
 B_5&0&J_1+J_3&D\\
 B_6&0&J_2&0\\
 B_7&K_5&0&0\\
 B_8&0&0&0.
\end{array}
\tag{D.2.22a16}
\]
All these operations are restrictions of the fiber tables,
so there is no unresolved Steenrod extension from the spectral
sequence.

The integral groups are
\[
\begin{gathered}
 H_5E=\ZZ/4,\quad H_6E=\ZZ/2,\quad H_7E=\ZZ/2,
 \quad H_8E=(\ZZ/2)^4,\\
 H_9E=\ZZ/8\oplus(\ZZ/2)^3.
\end{gathered}
\tag{D.2.22a17}
\]
The lower groups follow from the displayed first Bocksteins and
the known second Bockstein into $u_F$.
For the last group, the second Bockstein page in degree nine
is the single span of $N$, mapping injectively to that for $F$.
An incoming second differential would contradict its survival
in $F$. The classes $B_5,B_6$ cannot be earlier boundaries,
since their images have nonzero third differential in $F$.
Naturality therefore gives both
$\beta_3B_5=\beta_3B_6=N$ in $E$.
The complete degree-ten first-Bockstein cycle basis is
\[
 Z_2,U_2,T_{10},P_r,B_2,B_5,B_6,B_8,B_1+B_3.
\]
Under the Moore-four quadratic projection these map respectively
to $q_0,0,q_0,q_0,0,h+h',h+h',0,0$.
Thus only $B_5,B_6$ have nonzero third Bockstein, exactly as
required for \textup{(D.2.22a17)} and the ten-cell relations below.

In this particular degree the integral inclusion $H_9E\to H_9F$
is also injective. An element in its kernel reduces to zero modulo
two, by \textup{(D.2.22a13)}, so it lies in $2\ZZ/8\{N\}$.
The image of $N$ has odd coefficient on the unique order-eight
summand in \textup{(D.2.22a8)}, and hence has exact order eight.
No nonzero element of that even subgroup can vanish. Since
$F\to T$ is split, the integral Hurewicz kernel computed using
$E$ is the kernel in the statement.

\medskip\noindent
\textbf{The actual homotopy group and a stable cell model.}
The cover used above has an actual double-loop realization.
Put $Z=\operatorname{fib}(Y\to K(\mathbb F _2,4))$.
Then $F=\Omega^2Z$, the space $Z$ is five-connected, and
$\pi_6Z=\ZZ/2$.  Double-looping its first Postnikov map
$Z\to K(\mathbb F _2,6)$ realizes the map
$F\to K(\mathbb F _2,4)$, whose fiber is $E$.
Thus $E$ is a four-connected double-loop space.
The unit $E\to\Omega\Sigma E$ has a homotopy retraction by
the loop-space counit; Freudenthal makes it nine-connected.
Since $\Sigma E$ is five-connected, this proves
\[
 \pi_9T\cong\pi_9F\cong\pi_9E
       \cong\pi_9\Sigma^\infty E.
 \tag{D.2.22t1}
\]
This argument does not assert that stabilization of $F$ preserves
its specified double-loop operations.

Work two-locally.  Choose a finite stable approximation to
$\Sigma^\infty E$ through dimension eleven before cancelling
odd-unit cell pairs, and retain its ten-skeleton.
Write $L_i$ for its stable skeleta.  This convention removes
extraneous ten-cells killed by unit eleven-cell incidences and
preserves the required $\pi_9$.
Let $M=M(\ZZ/4,5)$, with bottom inclusion $i_5$ and top
six-class $D$, and put $r=D+E_6$ and $t=i_5\eta_5$.
The lower integral boundaries and operations in the preceding
tables give
\[
 \begin{split}
 D_0&=(M\vee S^6_r)\mathop{\cup}_{t+2r}e^7_S,
       \qquad L_7=D_0\vee S^7_R,\\
 D_0&=\cofib\left(
 \begin{pmatrix}4&\eta\\0&2\end{pmatrix}:
 S^5\vee S^6\longrightarrow S^5\vee S^6\right).
 \end{split}
 \tag{D.2.22t2}
\]
In particular $\beta S=r$ and $\operatorname{Sq}_*^2S=u_5$.
All cell assertions here concern stable cofibers; no minimal
unstable CW decomposition is being assumed.

Let $b\in\pi_7M$ be the order-two lift of $\eta$ on the
Moore top, and let $c=r\eta_6$.
The cofiber exact sequence gives
\[
 \pi_7L_7=\ZZ_{(2)}\{R\}\oplus
                    \ZZ/2\{b,c\}.
\]
The complete eight-cell attachment list, ordered as
$W,D_8,J_1,J_2,J_3$, is
\[
 0,\qquad b+c,\qquad b,\qquad0,\qquad2R+b.
 \tag{D.2.22t3}
\]
The $b,c$ coordinates are detected independently by the
second-square vectors $D,r$; the integral coordinate is fixed
by $\beta J_3=R$.
Subtract the $J_1$ cell from each of $D_8,J_3$.
These stable cofiber row changes give an actual splitting
\[
 \begin{split}
 L_8&=K_8\vee M(\ZZ/2,7)_R\vee S^8_W\vee S^8_{J_2},\\
 K_8&=D_0\cup_b e^8_{B_8}\cup_c e^8_{C_8},
       \qquad B_8=J_1,\quad C_8=D_8-J_1.
 \end{split}
 \tag{D.2.22t4}
\]
The subscripts $B_8,C_8$ here distinguish these cells from
the original space $B$ and from the cover's ten-classes $B_i$.

\medskip\noindent
\textbf{All ninth homotopy components of the eight-skeleton.}
First, $\pi_9M=\ZZ/4\{m\}$, with top quotient
$q_*m=2\nu_6$.  Write $z=r\nu_6$.
The full matrix cofiber sequence identifies the quotient of
$\pi_9D_0$ by $\ZZ/2\{z\}$ with the kernel of
\[
 \ZZ/8\{\nu_5\}\oplus\ZZ/2\{\eta_6^2\}
       \longrightarrow\ZZ/8,
       \qquad (x,y)\longmapsto4x+4y.
 \tag{D.2.22t5}
\]
This kernel is cyclic of order eight generated by
$(\nu_5,\eta_6^2)$; choose a lift $s_0$.
The possible discrepancy in $2s_0=m$ is $z$.
It is zero: the fourth-square detector of both $2s_0$ and
$m$ is zero, whereas the detector of $z$ is $r$.
The same argument excludes $8s_0=z$.
Changing $s_0$ by $z$ if necessary normalizes its
fourth-square vector to $D$.  Consequently
\[
 \pi_9D_0=\ZZ/8\{s_0\}\oplus\ZZ/2\{z\},
       \qquad2s_0=m.
 \tag{D.2.22t6}
\]
Here a square detector means the homology square on the top
class of the stable mapping cone.  In this range it is additive
and vanishes on twice an attaching map: the map from the cone
of $2f$ to the cone of $f$ is zero on the top mod-two class
and is the identity on the relevant lower homology.
The nonzero detector for $z$ is the Hopf-invariant-one
detector of $\nu$.

The two eight-attachments have composites
\[
 b\eta_7^2=2m=4s_0,
       \qquad c\eta_7^2=4z=0.
\]
There is no additional $\eta$-lift in their cofiber sequence:
the classes $b\eta_7,c\eta_7$ are independent before those
eight-cells are attached.  Therefore
\[
 \begin{split}
 \pi_9K_8&=\ZZ/4\{s\}\oplus\ZZ/2\{z\},\\
 \pi_9L_8&=\ZZ/4\{s\}\oplus\ZZ/2\{z\}
     \oplus\ZZ/4\{\tau\}
     \oplus\ZZ/2\{w\eta_8,j\eta_8\}.
 \end{split}
 \tag{D.2.22t7}
\]
The class $\tau$ is the Moore $\eta$ coextension in
$M(\ZZ/2,7)_R$, with $2\tau=R\eta_7^2$;
$w,j$ are the two free eight-sphere inclusions.
The stable stem and Moore identities used in these exact
sequences are the opened formulas
$\pi_1^S=\ZZ/2\{\eta\}$,
$\pi_2^S=\ZZ/2\{\eta^2\}$,
$\pi_3^S{}_{(2)}=\ZZ/8\{\nu\}$,
$\eta^3=4\nu$ and $2\widetilde\eta=i\eta^2$ in
\cite[\S2, (2.1)--(2.2), p.~149]{I_Oka1986}.

\medskip\noindent
\textbf{The exact $\eta$ composites of the free lifts.}
Set $J=M\cup_b e^8_{B_8}$.
Its cofiber sequence gives $\pi_9J=\ZZ/2\{\overline m\}$,
and $J/S^5$ is $C\eta$ with bottom cell in dimension six.
Let $e\in\pi_8J$ have top-cell degree two on $B_8$.
Its quotient is the top-degree-two class of $C\eta$.
The actual identity
$\eta_{C\eta}=i\nu q$, opened in
\cite[(8.4), p.~1279]{EH5AngeltveitRognes2005}, gives
\[
 e\eta_8=\overline m\quad\hbox{in }\pi_9J,
       \qquad e\eta_8=2s\quad\hbox{in }\pi_9L_8.
 \tag{D.2.22t8}
\]
Indeed its quotient is $2i\nu$, the nonzero image of
$\overline m$.  The free lift of $2C_8$ lies in the actual
$C\eta$ subcomplex on $r,C_8$ and has $\eta$ composite
$2r\nu=2z=0$.  The other two free composites are
$w\eta_8,j\eta_8$.

The degree-eight torsion of $L_8$ is
$\ZZ/4\{n=i_5\nu_5\}\oplus\ZZ/2\{\ell=R\eta_7\}$;
its free Hurewicz lattice has basis
$W,J_2,2J_1,2(D_8-J_1)$.
The nine-cells supply all four integral boundary directions.
The zero-boundary combinations $N$ and
$K_\Sigma=K_2+K_3+K_4$ kill $n,\ell$.
On taking $\eta$ composites, $n\eta=0$,
$\ell\eta=2\tau$, and \textup{(D.2.22t8)} gives $2s$.
Torsion additions to the integral attachments contribute only
$2\tau$, already killed.  Thus the image of
\textup{(D.2.22t7)} in $\pi_9L_9$ is exactly
\[
 V=(\ZZ/2)^5\{s,z,\tau,w\eta_8,j\eta_8\}.
 \tag{D.2.22t9}
\]
The quotient $\pi_9L_9/V$ is free of rank four, the kernel
of the full integral nine-cell attaching matrix.  In particular
$V$ is the entire torsion subgroup.  Its complete detectors are
\[
\begin{array}{c|c|c}
 &\operatorname{Sq}_*^2&\operatorname{Sq}_*^4\\ \hline
 s&0&D\\
 z&0&r\\
 \tau&J_1+J_3&0\\
 w\eta_8&W&0\\
 j\eta_8&J_2&0.
\end{array}
 \tag{D.2.22t10}
\]
The five vectors are independent.  Thus specifying these
coordinates leaves no unmeasured torsion component.

\medskip\noindent
\textbf{The free lattice, retaining the even-$\nu$ coefficients.}
Replace $N$ by $N'=N-K_\Sigma$, rescale $n$ by an odd
unit, and change $\ell$ by an even multiple of $n$.
We may then take the attachments of $N'$ and
$Q=K_\Sigma$ to be $n$ and $\ell$, respectively.
The order-eight homology generator changes only by an
order-two class.  The other two zero-boundary nine-cycles
have attachments
\[
 K_1:\ 2a n,
       \qquad K_5:\ 2b n,
       \qquad a,b\in\{0,1\}.
 \tag{D.2.22t11}
\]
These two bits are retained.  A mod-two spherical representative
alone would not set either bit to zero: its integral Hurewicz
class could differ from $K_i$ by $2N'$.

Choose integral order-two representatives $K_1,K_5,Q$ by
the half-boundaries of $B_3,B_7,B_4$, respectively, and use
the order-eight generator $N'$ supplied by the integral
configuration-chain calculation above.
The free Hurewicz lattice of $\pi_9L_9$ then has a basis
$p,q,x,y$ satisfying
\[
 \begin{aligned}
 h(p)&=4N',& h(q)&=2Q,\\
 h(x)&=K_1-2aN',&h(y)&=K_5-2bN'.
 \end{aligned}
 \tag{D.2.22t12}
\]
Normalize $p,q$ so that their second- and fourth-square
detectors are zero.  To justify this normalization, the Adem
relation $\operatorname{Sq}^2\operatorname{Sq}^2
=\operatorname{Sq}^3\operatorname{Sq}^1$ forces either
second-square vector into
$\langle W,J_2,J_1+J_3\rangle$.
For $p$ the top third square is zero; for $q$ it is $R$,
whose Bockstein is zero.  Either fourth-square vector lies
in $\langle D,r\rangle$.
The five entries of \textup{(D.2.22t10)} remove precisely
these five coordinates.  Even multiples of all attaching maps,
including $2x,2y$, have zero primary detectors by the
degree-two mapping-cone argument above.

\medskip\noindent
\textbf{All twelve ten-cell relations.}
The exact integral normalization is controlled by
$H_9E=\ZZ/8\{N'\}\oplus\ZZ/2\{K_1,K_5,Q\}$.
A first-Bockstein cycle with zero third Bockstein can be
made an exact zero-boundary ten-cell: its remaining $N'$
coefficient is divisible by sixteen and is removed by twice
an eight-$N'$ cell, while its order-two coefficients are
removed by even multiples of the corresponding degree-two
cells.  These changes preserve all primary detectors.
Since \emph{both} $B_5$ and $B_6$ have third Bockstein $N$,
they may each be normalized to boundary $8N'$.
Odd units do not change the ensuing two-primary relations.

The zero-boundary cells $Z_2,B_2,T_{10}$ give
\[
 w\eta_8=0,
       \qquad j\eta_8=0,
       \qquad w\eta_8+s+z=0.
 \tag{D.2.22t13}
\]
The two eight-$N'$ cells, using the normalized $p$, give
\[
 B_5:\ 2p+\tau+s=0,
       \qquad B_6:\ 2p+j\eta_8=0.
 \tag{D.2.22t14}
\]
Consequently $2p=0$ and $\tau=z=s$.
The zero-boundary cycle $B_1+B_3$ repeats $j\eta_8$.
The classes $U_2,P_r,B_8$ have zero detectors and impose
no new relation, by completeness of \textup{(D.2.22t10)}.
The third-Bockstein-zero combination $B_5+B_6$ repeats
$\tau+s+j\eta_8$.  Exactly one nonzero element of $V$
survives, represented by $s$.

The remaining integral cells have boundaries
$2Q,2K_1,2K_5$, which in \textup{(D.2.22t12)} are
$q,2x+ap,2y+bp$.  Their complete relations are
\[
 \begin{array}{ll}
 B_4:&q+\tau+s+j\eta_8=0,\\
 B_3:&2x+ap=0,\\
 B_7:&2y+bp=0.
 \end{array}
 \tag{D.2.22t15}
\]
Here $q$ was normalized above, and $2x,2y$ have zero
primary detectors; $B_1$ differs from $B_3$ by the already
killed $j\eta_8$.  All twelve ten-cells and all independent
integral boundary directions have now been used.
The resulting full presentation is
\[
 \pi_9E=
 \langle p,x,y,s\mid
       2p=2s=0,\quad2x=ap,\quad2y=bp\rangle.
 \tag{D.2.22t16}
\]
In particular, independently of the two undetermined bits,
\[
 4\pi_9T=0,
 \qquad
 \ker\bigl(h_{\ZZ}:\pi_9E\to H_9E\bigr)
       =\ZZ/2\{s\}\ne0.
 \tag{D.2.22t17}
\]
Its Hurewicz image is given exactly by
\textup{(D.2.22t12)} and has order eight.
The group is $(\ZZ/2)^4$ if $a=b=0$, and
$\ZZ/4\oplus(\ZZ/2)^2$ otherwise.
These alternatives are retained rather than selected by
an abstract group-order argument.
The previously proved injectivity of $H_9E\to H_9F$
and the splitting $F\to T$ identify the same kernel for
the integral Hurewicz map on $\pi_9T$.

Finally the $r$ sphere maps to an odd multiple of the
class $\rho\in\pi_6T=\ZZ/4$ of D.2.18.
Since $z=r\nu_6$ and $z=s$ in the final group, the
surviving element has the specified actual representative
\[
 s=\rho\nu_6\ne0.
 \tag{D.2.22t18}
\]
The same eight- and nine-cell exact sequences kill $n,\ell$
and the free directions $2J_1,2(D_8-J_1)$, while leaving
the $W,J_2$ directions modulo two.
Thus the actual group $\pi_8E=\pi_8T\cong(\ZZ/2)^2$ is
represented by those two eight-spheres.  Their $\eta$
composites are precisely the two classes killed in
\textup{(D.2.22t13)}.  Hence
\[
 \pi_8T\circ\eta_8=0.
 \tag{D.2.22t19}
\]
This excludes that particular proposed change of a relative
eight-cell filling.  It does not exhaust every possible change
of the preceding map and does not evaluate
$\lambda(\chi,\chi)$.


\medskip\noindent
\textbf{The common attachment coefficient and the remaining scope.}
The component calculation retains bits $a,b$ for the even-nu
parts of the $K_1,K_5$ attachments. They are equal, although
their common value is not determined here.
Indeed the actual homotopy Browder bracket of the bottom
$a\in\pi_2T$, of order two, with $\rho\in\pi_6T$ has order
dividing two. Its mod-two Hurewicz image is
\[
 \{a,b^2+k+p_3\}=K_1+K_5.
\]
The derivation law gives $\{a,b^2\}=0$ and
$\{a,p_3\}=\{a,f\}+aR=K_1$.
In the component presentation an integral Hurewicz class reducing
to $K_1+K_5$ has order dividing two only when $a=b$;
otherwise its coefficient on the order-eight class is $2$ modulo
$4$. Thus only the common bit remains. Neither choice changes
\textup{(D.2.22a1)--(D.2.22a2)}.
No identification of the actual square in \textup{D.2.21} with
the nonzero Hurewicz-kernel generator has been made. Its integral
Hurewicz image can nevertheless be determined exactly.

In integral homology the fiber generator is
\[
 h_{\ZZ}(\chi)=a^2-2c\in H_4T,
 \tag{D.2.22a19}
\]
not the uncorrected Pontryagin square $a^2$.
To verify the normalization, use the homomorphism
$K(\ZZ,2)=\mathbb CP^\infty\to K(\ZZ/2,2)=B$.
Its degree-four image is an odd generator $c$ because the
mod-two square of the fundamental class pulls back nontrivially.
The degree-two generator of $H_*\mathbb CP^\infty$ has
Pontryagin square twice the degree-four generator, as follows
from its primitive integral cohomology generator.
Consequently $q_*(a^2)=2c$. The kernel in degree four is
the order-two class in \textup{(D.2.22a19)}, which reduces to
$A$ and is the Hurewicz image of the bottom sphere of $F$.

The integral Poisson derivation law now gives
\[
 \{a^2-2c,a^2-2c\}=4\{c,c\}=8\xi c=0.
 \tag{D.2.22a20}
\]
Every omitted cross term has a factor $2a=0$; the even-label
self-bracket is $\{c,c\}=2\xi c$, and the exact order-eight
calculation above proves the last equality.
Naturality of Hurewicz for the homotopy Browder operation
therefore places the actual square in the kernel already
computed, proving \textup{(D.2.22a18)}. Since $4\pi_9T=0$,
four-divisibility in \textup{D.2.21} is exactly vanishing.
This does not distinguish the two values in
\textup{(D.2.22a18)}. The later geometric arguments settle
both retained questions: \textup{D.2.23} proves that the actual
square is zero, and \textup{D.2.31} proves that the common
attachment bit is zero, so $\pi_9T\cong(\ZZ/2)^4$.
The two-case calculation above is retained to identify exactly
which information those later constructions supply.
\end{proof}
